% ============================================================
\chapter{Implementation}
\label{chap:design}
% ============================================================

\section{Node Discovery and the Address Book}
\label{sec:discovery}

Node discovery merges addresses from three non-blocking sources into a single table.

\textbf{UDP probes.} The daemon opens two sockets on UDP port \texttt{4242}, one
joined to multicast group \texttt{239.255.0.99} and one bound for broadcast,
both with address and port reuse. The probe is a fixed magic string and the
reply carries a JSON node descriptor; it is a bare request-response probe, not
mDNS. The multicast socket exists because a UDP broadcast to a shared port is
delivered by the kernel to only one socket, so two daemons on one host would
never see each other; multicast with port reuse reaches all of them.

\textbf{TCP fallback.} Access points with client isolation drop broadcast
traffic without error. When the UDP probe returns nothing, the daemon sweeps
each local interface's prefix (container and VPN interfaces skipped) with
concurrent health probes, 64 at a time with a 700\,ms timeout each, and stops
once any daemon answers. A reply's host is taken from the packet's source
address, never from the address the peer advertised at startup.

\textbf{The address book.} The node store keeps every node under its node id,
the fingerprint of its key, together with every \emph{path} to it: its kind
(LAN, internet, gossip, self), host, port, and when it was last seen, last
worked and last failed. A sighting only adds a candidate. A path becomes the
active one only after the poller has reached it and the certificate's key
matched the node id; an active path that still answers is kept, and a verified
LAN path beats a verified internet path. Members also gossip their own
addresses, signed with their node keys, over the authenticated API; a gossiped
address is a candidate, never a fact.

The poller runs every 10\,s, trying the active path first; a node with no
verified path for 15 minutes is removed, and a sighting alone never refreshes
it.

% ------------------------------------------------------------
\section{Joining Across Networks}
\label{sec:nat}

A broadcast cannot leave the local network, and a machine behind another router
accepts no inbound connection. Figure~\ref{fig:nat} shows how two such machines
come to share a direct, encrypted link without either opening a port by hand and
without any machine relaying their traffic.

\begin{figure}[!htbp]
    \centering
    \includegraphics[width=124mm]{\figpath{fig_4.1_nat_traversal.png}}
    \caption[Joining across networks]{Joining across networks. The introducer exchanges candidate
    addresses and forwards one punch request; the two nodes then race their
    candidate pairs over a single UDP socket, fall back to a birthday collision
    when one NAT allocates a port per destination, and carry everything over a
    mutually authenticated QUIC link exposed locally as an ordinary node.}
    \label{fig:nat}
\end{figure}

\textbf{One command per machine.} \texttt{pipedpeer invite} prints a single
line, \texttt{pipedpeer connect <code>}, whose code carries the cluster secret
and up to six introducer addresses with a checksum, so a mistyped code is
refused as damaged rather than tried. \texttt{connect} saves both, restarts the
daemon once, waits up to a minute, and ends with one of three results:
connected, naming the machines (exit 0); partly connected, giving each missing
machine's reason and fix (exit 2); or failed, with what to do (exit 1).

\textbf{An introducer in every daemon.} Every daemon holding a secret serves the
introducer for its own cluster on UDP 4244 and refuses other clusters. A node
registers with every introducer it knows --- those in its invite code and
members it has verified, at most eight --- every 30\,s and at once when its own
address changes, so the cluster survives the loss or readdressing of any one
member. A public server can join as an introducer that takes no jobs. Nodes
register their candidates under a cluster identifier derived from the secret by
hashing, so an introducer never holds the secret itself. A \texttt{connect}
request makes it forward a single punch datagram to the target, carrying only
the candidates the caller itself registered, rate-limited to four requests per
second per source. It carries no job data.

\textbf{Candidates.} Each node offers three kinds of address: its LAN addresses;
a port its router agreed to forward, requested by PCP \cite{rfc6887}, NAT-PMP
\cite{rfc6886} or UPnP-IGD \cite{upnpigd} for two hours and renewed at half the
lifetime granted; and the reflexive address the introducer observed, in the
manner of STUN \cite{rfc8489}. A router mapping is withheld when its external
address differs from the reflexive one: behind a second layer of NAT, such as a
phone tether, the router that answers is not the one facing the internet.

\textbf{One socket.} Punch probes, introducer messages and QUIC share a single UDP
socket on port 4243. Sharing one socket is required for correctness, not convenience: a router's
mapping belongs to the socket that created it \cite{ford2005p2p}, so the
connection must arrive on the socket the probes left from. Datagrams are
separated by content --- probes by four magic bytes, introducer messages by a
leading brace, everything else handed to QUIC.

\textbf{The race.} Both sides probe every pairing of their candidates every
250\,ms for fifteen seconds, rather than trying candidates in turn, which would
add a timeout for every address that cannot work. When one side's NAT allocates
a fresh external port for every destination there is no address to publish, and
the \textbf{birthday collision} \cite{anderson2020nat} runs: that side opens 256
sockets toward the peer, the other sends 600 probes to random ports among the
64{,}511 above 1024, and the socket that is hit is kept, since its mapping is
the reachable one. A round hits one of the 256 mappings with probability
$1-(1-256/64511)^{600} \approx 91\%$.

\textbf{The link.} The winning path carries QUIC \cite{rfc9000} with TLS~1.3.
Inside it each side sends a hello signed with its Ed25519 node key, and each
checks that the key's fingerprint is the node it meant to reach and that the
cluster identifier matches. The hello signs the session's exported keying
material \cite{rfc5705} and proves the cluster secret with an HMAC
\cite{rfc2104} over the same value, so it is valid for one session and one
cluster only and cannot be replayed to a third node. The link is exposed to the
rest of the daemon as a TCP forwarder on a loopback port, recorded as that
node's internet path, and each local connection becomes one QUIC stream to the
peer's daemon. Placement, closure transfer and gradient exchange therefore work
over the link unchanged.

\textbf{Roaming.} A hello from a key that already has a link replaces that link
at once. The daemon checks its interfaces every 5\,s; on a change it migrates
the QUIC connections it dialled to the new address, redials those it accepted,
and registers again with every introducer.

\textbf{No relay.} A relay carries every byte of every job and becomes the
cluster's bottleneck and its cost, so none is in the data path. A peer with no
direct path is logged with the reason, never offered to the scheduler, and
retried after a back-off; \texttt{pipedpeer status} lists it with that reason.

% ------------------------------------------------------------
\section{Security and Identity}
\label{sec:trust}

Security has three layers (Figure~\ref{fig:trust}).

\begin{figure}[!htbp]
    \centering
    \includegraphics[width=144mm]{\figpath{fig_4.2_trust.png}}
    \caption[The three layers of trust]{The three layers of trust: a shared secret on every API route,
    pinned TLS (or QUIC across networks) on the wire, and the node key that
    makes the certificate, authenticates links and signs closures.}
    \label{fig:trust}
\end{figure}

\textbf{The shared secret.} The first machine's setup stores 32 random bytes,
readable only by the user, before its daemon first starts; \texttt{connect}
stores the secret from the invite code. Every route on the daemon,
\texttt{/health} included, requires it, compared in constant time, and outbound
requests carry it automatically. Loopback is not exempt, since peers across
networks arrive through loopback forwarders. The secret is attached only to a
host that has proved it holds it, by returning an HMAC over a challenge. Jobs
receive it so that the shim can reach its own daemon, and the web page receives
it in the URL fragment opened by \texttt{pipedpeer ui}, which a browser never
sends to a server.

\textbf{Pinned TLS.} Each daemon's self-signed certificate is made from its
Ed25519 node key. TLS and plain HTTP share the daemon port; the listener inspects
the first byte of each connection. Clients try TLS first, and the pin is keyed
by node id rather than by address: a certificate whose key matches the expected
node id is accepted wherever the node now is, and a different key for a known
node id is refused. The scheme is trust on first use and opportunistic: a
handshake that fails for a reason other than a pin mismatch makes the client
remember that peer as plain HTTP. Traffic across networks is encrypted by QUIC.

\textbf{Node keys and signed closures.} Each node holds one Ed25519 key
\cite{bernstein2012ed25519} whose fingerprint is its node id everywhere. The same
key authenticates QUIC links, makes the TLS certificate, signs gossiped
addresses, and signs closures as a Nix signing key. Nodes publish the public half
in their heartbeat, each daemon keeps a list of the keys of the nodes it knows,
and \texttt{setup}, with consent, adds them to the system Nix configuration as
trusted keys; it never disables Nix's signature requirement. A closure sent
between daemons in the signed format is imported under the receiver's own
signature policy.

A worker learns a new peer's key before Nix does, so the import passes the
cluster's keys to that one copy operation, after checking the archive's named
signer against the signer's published identity. Because a multi-user Nix daemon
ignores keys passed by an untrusted user, setup also installs a root-owned
helper that can do nothing but add or remove keys of the cluster's fixed form.
The limits of this design are set out in Section~\ref{sec:security}.

% ------------------------------------------------------------
\section{Placement and Scoring}
\label{sec:placement}

Candidates are merged from the address book, live discovery and the local node,
deduplicated by node id; a node with a role is never offered a job. Five hard
filters run in order: architecture must match the closure; a GPU must be present
if one was required; and free memory, free cores and free VRAM must each satisfy
the request (Figure~\ref{fig:placement}). Every rejection records a reason, shown by
\texttt{pipedpeer jobs <id> -{}-explain}.

\begin{figure}[!htbp]
    \centering
    \includegraphics[width=150mm]{\figpath{fig_4.3_placement_flow.png}}
    \caption[Placement flow]{Placement: gather, filter, score, lease, retry. A task is never
    dropped; when nothing is eligible it is queued and retried.}
    \label{fig:placement}
\end{figure}

Survivors are scored. Writing $u_{\mathrm{cpu}}$ and $u_{\mathrm{mem}}$ for CPU and
memory utilisation in percent, $n_{\mathrm{jobs}}$ for the count of running jobs,
$\hat{c}, \hat{f}, \hat{m}, \hat{r}$ for the core count, clock speed, free
memory and free-core ratio each normalised onto $[0,1]$, and $h \in [0,1]$ for
the node's health score, the score is
\begin{equation}
\label{eq:score}
\begin{split}
S = \Bigl(1 &- \frac{u_{\mathrm{cpu}} + u_{\mathrm{mem}}}{200}
     - 0.05\,n_{\mathrm{jobs}} + 0.10(\hat{c}-0.5) \\
    &+ 0.06(\hat{f}-0.5) + 0.10(\hat{m}-0.5) + 0.10(\hat{r}-0.5)\Bigr)\cdot h .
\end{split}
\end{equation}
Two properties matter more than the individual coefficients. \textbf{Load
dominates hardware}: the two utilisation terms can subtract $1.0$, while the
four hardware terms together move the score by at most $\pm0.18$, so a fast but
busy machine loses to a slower idle one, as it should on
workstations that people are also using. \textbf{Missing telemetry is neutral}:
normalisation returns the midpoint $0.5$ for an absent value, so a node that does
not report its clock speed is neither rewarded nor punished. $S$ is clamped at
zero.

When a GPU is requested, a second function adds up to $0.8$ to $S$, drawn from total
VRAM, device count, free VRAM on the best device and compute capability, less a
utilisation penalty. The free memory a node advertises is what it can actually
\emph{grant}: its own budget less the reservations of every daemon on the same
machine (Section~\ref{sec:sandbox}).

A node that refuses the job's signed closure is left out for the rest of the
run, and the client prints the command that fixes it on that node.

% ------------------------------------------------------------
\section{The Lease Protocol}
\label{sec:lease}

Placement is a three-phase handshake, and its state machine is where crash
recovery lives (Figure~\ref{fig:lease}).

\begin{figure}[!htbp]
    \centering
    \includegraphics[width=150mm]{\figpath{fig_4.5_lease_fsm.png}}
    \caption[Lease states]{Lease states. Recovery is a consequence of expiry rather than a
    separate mechanism.}
    \label{fig:lease}
\end{figure}

An \texttt{accept} reserves capacity for 30 seconds with a 5 second grace; a
\texttt{commit} promotes the reservation to running, renewed every 30 seconds
and reaped after 5 minutes without renewal; \texttt{complete} releases it. A
sweeper runs every 2 seconds. The submitter renews while the environment builds
and uploads, so a long build never outlives its lease, and a lease bound to a
running job is kept past its expiry.

Two details carry the correctness argument. The concurrency and memory checks
and the lease insertion happen \textbf{under a single lock}, so two submitters
racing for the last slot cannot both succeed. And the concurrency limit counts
reserved and running leases together, so a burst of submitters cannot slip
through the window between accept and commit.

Recovery then needs no further machinery. A submitter that dies stops renewing,
and the sweeper reclaims its slot. A worker that dies stops answering health
checks and its in-flight connection fails, which the coordinator treats as a
failed attempt and reschedules, with back-off doubling from 2 seconds to a
2 minute ceiling. Waiting for room never drops a task, while a job that
\emph{fails} is run at most three times, on untried nodes first. From the
submitter a dead node and a lost network look the same, but the first means run
the job elsewhere and the second means wait; the client therefore runs three
probes at once, two seconds each and none needing root --- another peer's health
endpoint, a registration with the introducer, and a TCP connection to the
default gateway --- and an answer to any of them confirms that the remote peer,
not the local network, failed.

% ------------------------------------------------------------
\section{Reproducing the Environment}
\label{sec:env}

\textbf{Import scanning.} The client extracts imports with a line-oriented
parser: it strips comments, takes the module from \texttt{from X import \ldots},
and splits \texttt{import a, b as c} into every module named. It does not see an
import performed through \texttt{importlib}. Each name is classified as a local
file, a standard-library module, or an external dependency, and external names
are resolved through a mapping table (\texttt{sklearn} to
\texttt{scikit-learn}, for example) and otherwise looked up as a package of the
same name in the pinned package set, with one evaluation for all of them. The
scan follows the script's own modules, with a visited set so that import cycles
end; a \texttt{requirements.txt} beside the script is read, and a
\texttt{uv.lock} file or an explicit \texttt{-{}-pkg} flag override detection.
Only what is still unresolved is reported, before anything runs, with the flag
that would add it.

\textbf{Content-addressed environment pinning.} Nix \cite{dolstra2004nix} names every build output by a
hash of its inputs. The generated flake pins an \emph{exact} nixpkgs revision,
not a branch. A branch resolves at build time, so two nodes building the same
script on different days would produce different store paths and the system
would re-ship a multi-hundred-megabyte archive for every job. Pinning makes the
store path a stable, content-derived name for an environment, identical across
machines and across time, and therefore a cluster-wide cache key
(Section~\ref{sec:res-cache}).

Every environment also contains \texttt{cloudpickle} \cite{cloudpickle}, so that
parallel kernels with no source can be shipped by value. On x86\_64 Linux,
scikit-learn is pinned to a specific wheel, and torch with CUDA is fetched as a
pinned wheel set inside a \emph{fixed-output derivation} with a recorded hash,
because building it from source takes hours. The CUDA build is used when a GPU
is requested, which is the default when a script imports torch.

\textbf{Common packages built ahead of time.} Environments share Nix store
paths, and a transfer sends only the paths a peer lacks. After
\texttt{connect}, every machine builds in the background one environment
importing the usual scientific set (numpy, pandas, scikit-learn, scipy,
matplotlib, seaborn, tqdm and CUDA torch), so a script seen for the first time
shares those paths and builds only what it adds.

% ------------------------------------------------------------
\section{Closure Transfer and Cache Reuse}
\label{sec:cache}

\begin{figure}[!htbp]
    \centering
    \includegraphics[width=150mm]{\figpath{fig_4.6_closure_cache.png}}
    \caption[Closure cache and transfer]{The store path is the cluster-wide cache key. Most submissions ship
    no environment at all, and a related environment ships only the store paths
    the receiver lacks. The upload never waits more than 60\,s for a slow peer.}
    \label{fig:closure}
\end{figure}

Before uploading, the client asks the target whether it can already run the
closure (Figure~\ref{fig:closure}), and only on a negative answer does anything cross the wire. Two
questions are needed: whether the archive is in the node's cache, and whether the
closure is \emph{materialised}, that is, whether its entry point exists on disk.
A deployment with a shared store volume has the closure present without ever
having cached an archive, while on a per-node store the entry point appears only
after an import; either predicate alone is wrong for one of the layouts.

\textbf{Path-level difference.} A closure is a set of store paths, and two
environments that differ by one package share most of them. The client lists the
closure's paths and asks the target which it lacks; if some but not all are
missing, only those are sent, as a partial import that the receiver never caches
as the whole closure. A fresh node may instead fetch public paths from the Nix
binary cache; if that download has not finished 30 seconds after the sender's
archive arrived, it is abandoned and the sender sends every missing path itself.

\textbf{Signed transfer between daemons.} After an upload, the worker pushes the
closure to peers that cannot yet run it. Parallel work is shared only with peers
that can run the closure, so this push is what makes the cluster usable for the
distribution of Section~\ref{sec:intercept}. A peer that advertises the signed
format receives one archive containing a signed binary cache --- the signature
records for every path and the compressed archives for the paths it lacks ---
produced and imported by Nix's own copy mechanism.

\textbf{Not waiting for the slowest peer.} The upload waits at most 60 seconds
for these pushes. A peer still importing finishes in the background and joins
the pool by itself, since each parallel chunk asks afresh which peers can run
it; meanwhile it reports ``importing'', so no second send starts. A background
send is cut only when it stops moving, after 60 seconds with no bytes sent.
Both closure caches are pruned least-recently-used to 8\,GiB each, and workspace
and results archives are compressed with Zstandard \cite{rfc8878}.

% ------------------------------------------------------------
\section{Execution and Isolation}
\label{sec:sandbox}

\begin{figure}[!htbp]
    \centering
    \includegraphics[width=130mm]{\figpath{fig_4.7_sandbox.png}}
    \caption[The sandbox]{What the job can and cannot see. Shared memory is mounted
    world-writable because Python's \texttt{multiprocessing} requires it, and
    the driver libraries are bound file by file so that the host C library
    cannot shadow the closure's.}
    \label{fig:sandbox}
\end{figure}

Jobs run under \texttt{crun} \cite{crun} in an OCI bundle \cite{oci2021} with a
minimal read-only root (Figure~\ref{fig:sandbox}): the store read-only, the workspace read-write, a private
home, and temporary filesystems for scratch space. Shared memory is mounted
world-writable because Python's \texttt{multiprocessing} maps POSIX shared memory
for its semaphores. GPU support mounts the device nodes and binds the driver
libraries \emph{file by file}, because binding the host library directory would
place the host C library ahead of the closure's, and torch would then fall back
to the CPU silently.

\textbf{A memory ceiling for every job.} Where the daemon has a usable cgroup v2
hierarchy, each job runs in its own cgroup with a memory limit, and an equal swap
limit, of twice its estimated requirement with a 512\,MiB floor, clamped to what
the node can grant but never below what admission approved; a job with no
estimate is capped at 4\,GiB. When a job is killed and the cgroup's
out-of-memory counter has risen, the client says so and how to raise
\texttt{-{}-mem}.

\textbf{A memory ceiling for the daemon.} Parallel chunks run in warm worker
processes in the daemon's own cgroup, so \texttt{pipedpeer start} launches the
daemon in a user \texttt{systemd} scope with a ceiling of half the machine and
swap refused. A daemon without a user \texttt{systemd}, such as one in a
container, runs a watchdog instead that measures the daemon's children every
second and, past the ceiling, kills the largest child and logs it, never the
daemon. A hard limit is used rather than a throttling one, because a job held
at a throttling limit stalls instead of failing. Daemons on a single machine,
recognised by the kernel's boot identifier, subtract one another's reservations,
so two of them cannot promise the same memory twice.

\textbf{Without root.} On a machine with no system Nix, setup seeds a private
store in the user's data directory and mounts it at \texttt{/nix} inside a user
and mount namespace, and installs \texttt{crun} into the user's own binary
directory when no package manager supplies it.

\textbf{What the sandbox does not do.} The job runs as user 0 inside the
container and shares the host's network; the specification sets no seccomp
filter and no capability list. The namespaces in use are process, IPC and mount,
plus hostname where the kernel allows it, and a user namespace when \texttt{crun}
runs without root. A node without \texttt{crun} or user namespaces warns and runs
the job unisolated, unless it uses a private Nix store, in which case it refuses.

\textbf{Results and output.} Only files created or modified during the run are
returned, with deletions carried in a manifest, so a job cannot overwrite the
submitter's own sources; any path that escapes the job's directory is refused on
both sides. Output is forwarded as each read of up to 4\,KB arrives, so a
progress bar on a remote node animates on the submitter's terminal.

\textbf{Stopping and detaching.} An interrupt or a closed terminal makes the
client send one cancel frame; the daemon signals the job's process group, kills
it two seconds later if it is still alive, and the cancel reaches remote runs,
parallel chunks on peers and every training rank. A connection that ends
\emph{without} a cancel frame never kills the job, and background mode builds on
this: the daemon logs every output frame, a detach closes the channel without a
cancel, and \texttt{pipedpeer attach} replays the log, follows live, and
downloads the results at the end. A detached job's directory is pruned 24 hours
after it ends, and its log is capped at 50\,MB.

\textbf{Phones.} An Android phone joins under Termux \cite{termux} as a worker for
parallel chunks, running them from a Nix environment inside a Debian under
\texttt{proot} \cite{proot} where setup can install one, and otherwise with
Termux's own Python. Closures go only to peers of the same architecture, so a
phone receives the job's package list instead, builds its own environment in the
background, and declines the job until it is ready.

% ------------------------------------------------------------
\section{Transparent Interception}
\label{sec:intercept}

\textbf{Delivery.} The shim is a Python package embedded in the binary. It is
appended to the workspace archive under a hidden directory, so the user's project
on disk is never modified, and the worker points the interpreter's module path at
that directory, where a \texttt{sitecustomize} module makes CPython import it
before the user's first line. Nothing is installed on any node. Nix ships its own
module of the same name, which adds the environment's package directories; the
shim performs that work first, since only the first module on the path is
imported (Figure~\ref{fig:shim}). \texttt{multiprocessing.Pool} and \texttt{ProcessPoolExecutor} are
patched at startup, while NumPy, pandas, torch and joblib are patched on their
first import through an import hook, so a job that never imports torch never pays
its 1.7\,s import (Section~\ref{sec:res-shim}).

\begin{figure}[!htbp]
    \centering
    \includegraphics[width=142mm]{\figpath{fig_4.8_shim_injection.png}}
    \caption[Shim injection]{How the shim reaches the interpreter and what it replaces.}
    \label{fig:shim}
\end{figure}

\subsection{Parallel Maps}

\texttt{map}, \texttt{starmap}, \texttt{imap} and \texttt{imap\_unordered} go to
the cluster (\texttt{imap} lazily, in batches), while \texttt{apply} and
\texttt{apply\_async} stay local. A small prefix of items runs locally first to measure
per-item cost, and a map whose remaining work is under half a second stays
local. The kernel reaches the peer in one of two ways. A module-level function,
including one defined in the script's \texttt{\_\_main\_\_}, travels as its
\emph{source}, with the helper functions it calls and a pickle of the module data
it reads. A lambda, closure, bound method, partial or decorated function travels
\emph{by value} through \texttt{cloudpickle}. Anything neither can carry is
counted as unshippable and runs locally. Every run can write a \textbf{receipt}
counting the items that ran remotely and locally, the parts that failed, and the
node each part ran on, so a claim about where work ran can be checked.

\textbf{Sharing work by measured speed.} On the daemon, a chunk is divided among
peers by a makespan solver. Each device is modelled as a setup time $o_i$ plus a
throughput $r_i$, fitted by least squares from the last eight parts it returned;
an unmeasured peer is assumed to be as fast as this node, scaled by its core
count, since assuming it slow would size its share below anything worth sending.
The solver finds the finishing time $T$ at which the devices together complete
the work, allocating $(T - o_i)\,r_i$ items to each, and drops any device whose
setup time alone exceeds $T$. On a peer, a part fans out across its cores.
Division therefore repeats at two levels, machines and then cores, and a part
already split off is never split across machines or forwarded again. A map too
large for the node that holds it is the one case of forwarding: it goes whole,
once, to a peer that can hold it, and that peer makes the split
(Section~\ref{sec:ft}).

\textbf{Slow peers.} A part still unfinished past its prediction counts as taking
at least that long, so a slow peer's estimate can only grow. Three rules act
before the fit has caught up: a split part waiting for its peer starts on another
peer with a free slot when that would finish sooner; a part running past its
prediction gets a copy on the fastest free peer, the first result winning and the
loser cancelled; and a chunk too small to split goes to the peer predicted to
finish first. A peer that lost its last three races stops being raced.

\textbf{Transfer.} Each node's chunks in flight follow its measured round trip,
and chunk size follows the measured transfer cost. Large functions and globals
cross once per peer, by hash; requests are compressed at a level chosen by the
measured link speed; and numerical buffers travel outside the serialised
request. With \texttt{-{}-offload}, every distributed piece goes only to other
machines, keeping the sender free for its user.

\begin{table}[!htbp]
\centering
\caption{Intercepted primitives and their distribution strategies}
\label{tab:intercept}
\renewcommand{\arraystretch}{1.2}
\rowcolors{2}{}{rowgray}
\begin{tabularx}{\textwidth}{L{5.0cm} X}
\toprule
\textbf{Primitive} & \textbf{Strategy} \\
\midrule
\texttt{multiprocessing.Pool}, \texttt{ProcessPoolExecutor.map} & Measure a few items locally, then race local against remote chunks shared by measured speed \\
pandas \texttt{groupby().agg()}, \texttt{merge}, \texttt{join} & Hash shuffle on the key; bucket zero local \\
pandas \texttt{read\_csv}, \texttt{read\_parquet} & Out-of-core reads in 64\,MB chunks cut at record boundaries \\
\texttt{numpy.matmul}, \texttt{dot}, \texttt{tensordot} (2-D) & Split into row blocks, concatenated at the origin \\
\texttt{numpy.linalg.svd}, \texttt{eig} & Whole matrix offloaded to one worker \\
\texttt{joblib.Parallel}, scikit-learn \texttt{n\_jobs} & Batches of delayed tasks, shared arrays sent once per chunk \\
torch \texttt{matmul}, \texttt{mm} & Row blocks, gated on size \\
torch training (\texttt{-{}-ddp}) & Gradients or weights averaged by the lead daemon, weighted by samples \\
\bottomrule
\end{tabularx}
\end{table}

Interception replaces module functions, so the operator form \texttt{A @ B} is
not intercepted for NumPy arrays. A \texttt{joblib} call made inside a remote
task runs inline on that worker. pandas operations are distributed only above
512\,MB unless forced, because their arithmetic intensity is low.

\textbf{Why the shuffle is exact.} Hash partitioning on the key is the standard
technique of MapReduce \cite{dean2004mapreduce}; PipedPeer uses only that path,
never a combiner. Rows are bucketed by a hash of the grouping or join key, so
equal keys land on the same node and every key is complete on exactly one node
(Figure~\ref{fig:shuffle}).
Each node's aggregation is therefore a \emph{complete} aggregation, the origin
combines the partials by concatenation, and \emph{any} aggregation specification
works, including those with no combiner form. Exactness is of values, not of row
order: a concatenation of buckets does not reproduce pandas' join ordering, so a
script that depends on the row order of a \texttt{merge} result can see a
different order.

\begin{figure}[!htbp]
    \centering
    \includegraphics[width=117mm]{\figpath{fig_4.9_hash_shuffle.png}}
    \caption[Hash shuffle]{Hash shuffle. Correctness follows from key locality, not from
    combiner algebra.}
    \label{fig:shuffle}
\end{figure}

\subsection{Data-Parallel Training}
\label{sec:ddp}

\textbf{Synchronisation over the daemon port.} Rather than let torch build its
rank-to-rank mesh \cite{li2020pytorchddp}, which needs direct connections on its
own ports between every pair of ranks, each rank posts its gradient to the lead
node's ordinary daemon port, which works wherever the daemon port works,
including over the links of Section~\ref{sec:nat}. The daemon \emph{reduces}:
gradients arrive as raw frames, are summed in double precision weighted by the
number of samples behind each, $\sum_i n_i g_i / \sum_i n_i$, and one averaged
gradient is returned (Figure~\ref{fig:ddp}). Returning every rank's gradient to every rank instead was
measured 21 times slower on two machines (62.9\,s against 2.9\,s). Gradients
travel as 8-bit integers with one scale per block of 4096 values, the rounding
error carried into the next step in both directions (error feedback)
\cite{dettmers2016eightbit, seide2014onebit}. Three or more ranks reduce in a
butterfly rather than through one hub, and ranks on one machine average locally
first.

\textbf{Choosing the ring.} Each candidate's throughput is probed, and the ring
admitted is the prefix of the fastest that maximises $k \cdot r_k$ --- the ring's
worth if every rank ran at the slowest admitted rate --- using one device kind
only, GPUs when any are wanted and present, with the lead on a GPU node. Each
rank's shard and batch slice are sized by its share, and the slices sum to the
script's own batch, so the learning rate and the mathematics are the script's.
Every rank sends a manifest of its parameters to the opening barrier, and a
mismatch stops every rank with one message. Every 20 synchronisations the daemon
refits the shares from observed step rates, applying them at the next epoch
boundary when a share has moved by a tenth or more, and it can drop the slowest
rank --- never rank 0, never below two --- when that rank adds more
synchronisation time than computation. A round that a straggler has not joined
completes without it after 45 seconds.

\textbf{One GPU or a split.} After a warm-up, the lead daemon compares the
measured step time of the split with that of the fastest rank alone, and either
keeps the ring or releases every rank but the fastest, which exit with a code
meaning ``released'' rather than ``failed''. A user who names a ring size with
\texttt{-{}-ddp~$N$} fixes the number of ranks, and the ring is then never
released or shrunk.

\textbf{Synchronising less often.} \texttt{-{}-ddp-sync fast} synchronises every
$H$ steps instead of every step, in the manner of DiLoCo
\cite{douillard2023diloco} and post-local SGD \cite{lin2020localsgd}: ten steps
of exact synchronisation as a warm-up, then local steps, with the ranks' weights
averaged by the number of samples each processed, every epoch ending on a
synchronisation, and the optimiser's moments averaged every tenth round. It never
adds epochs or changes the batch or the learning rate. The outer step is plain
averaging (outer learning rate 1, no momentum). DiLoCo's Nesterov outer step,
momentum 0.9 at an outer learning rate of 0.7, settles at $0.7/(1-0.9) = 7$ times
the averaged change; at intervals of a few hundred steps that is harmless, but at
the 2 to 64 steps used here it multiplies the script's learning rate, and it
failed the accuracy bar at every interval tried. The bar is one-sided: the final
validation loss may be at most 1\% above one device's, as a mean over three
seeds, on every test model. The interval depends on the ring's size: the largest
that passed on every model was 32 steps at two ranks, 16 at three and 4 at four,
and none passed at eight ranks, so none is used for five or more
(Section~\ref{sec:res-train}). That table is the
default; an explicit fast ring beyond four ranks runs at $H=4$ with a warning.

\textbf{Choosing automatically.} By default the lead daemon switches a ring to
fast synchronisation once, when every rank has reported, the ring is within the
validated size, and the predicted fast step, inflated by 10\%, beats the exact
split and, for a ring not fixed by \texttt{-{}-ddp~$N$}, also the fastest rank
alone. Under \texttt{-{}-ddp~$N$} the ring size is fixed and the automatic
choice is between the two sync modes only. Every rank switches at the same step boundary, rank 0
prints why, and a joining rank adopts the ring's mode; a join that takes an
automatic fast ring past the validated size returns it to exact synchronisation
at the next epoch boundary. A rank whose local round runs past four times the
ring's median step time skips that round's remaining batches but counts them, so
every rank stays on the same step.

\textbf{Losing and regaining ranks.} When a machine is lost, the client names the
rank and host, the lead daemon stops waiting for it, and the survivors take new
shares at the next epoch boundary. Samples are assigned to ranks by a weighted
rendezvous hash \cite{thaler1998rendezvous}, which keeps each survivor's share
exact and over 95\% of its samples in place, where hash-modulo routing would move
most of them. A rank whose script fails stops every rank at once. Losing rank 0's
machine ends the run, since its daemon holds the synchronisation, and the
survivors stop within about 30 seconds with that reason. A returning machine or a
new eligible node joins the ring at the next epoch boundary: it loads the lead's
model, optimiser and buffer snapshot through the lead daemon, runs one inert
batch for each epoch it missed so that the script's loop variables exist, and the
world size and shards grow to include it, moving about $1/n$ of the data. Only
scripts that use a data loader can be re-spread.

The first synchronisation is a barrier: every rank must reach the opening weight
broadcast within a configurable wait, five minutes by default, and ranks still
missing are dropped and named. The torch gloo and NCCL back-ends are not used.

\begin{figure}[!htbp]
    \centering
    \includegraphics[width=123mm]{\figpath{fig_4.10_ddp.png}}
    \caption[Training synchronisation over the daemon port]{Synchronisation over the daemon port. The lead daemon averages by
    samples and returns one average: int8 gradients every step in exact mode,
    weights every $H$ steps in fast mode; auto switches to fast once.}
    \label{fig:ddp}
\end{figure}

% ------------------------------------------------------------
\section{The Cost Model}
\label{sec:cost}

Interception is worthwhile only when the cluster wins, so numerical work is
gated. Writing $b$ for the payload in bytes, $\phi$ for its arithmetic intensity
in floating-point operations per byte, $\rho$ for the number of times the data
crosses the link, $B$ for link bandwidth and $K$ for the number of participating
nodes, divisible work is distributed only when

\begin{equation}
\label{eq:cost}
\underbrace{\frac{b\phi}{R}}_{\text{compute locally}}
\;>\;
\underbrace{\frac{\rho\, b}{B}}_{\text{move it}}
\;+\;
\underbrace{\frac{b\phi}{R' K}\times 1.3}_{\text{compute remotely}}
\end{equation}

where $R$ is the local throughput, $R'$ the remote throughput, and the factor
$1.3$ charges 30\% overhead against ideal parallel speedup. Matrix
multiplication is priced at about 200\,GFLOP/s with $\rho = 5$, pandas work at
1\,GFLOP/s. Work that cannot be divided --- an SVD or eigendecomposition, priced
at about 1.5\,GFLOP/s and sent whole to one worker --- moves when the transfer
costs less than half the local computation, since it frees the origin entirely.

The checks run in order. At least two nodes and 32\,MB of payload are required.
\textbf{If the working set does not fit in local memory, the work goes to the
cluster} whatever the transfer costs, and a matrix product the cluster cannot
take either raises an error naming the size and the remedy rather than exhausting
the machine. If any peer reports a GPU, $R'$ is taken as six times $R$, a conservative
heuristic multiplier rather than a measured speed. The inequality is
then evaluated at an optimistic 10\,Gbit/s (Figure~\ref{fig:cost}); if offloading loses even there, the
work stays local without measuring anything. Otherwise bandwidth is
\emph{measured}: the daemon's rate over recent real transfers is used when it
has one, and otherwise an 8\,MB probe is pushed through the real dispatch path
(about 3\,s on a 2.5\,MB/s link); the result is cached for five minutes, and a
failed probe keeps the work local.

\begin{figure}[!htbp]
    \centering
    \includegraphics[width=80mm]{\figpath{fig_4.11_cost_model.png}}
    \caption[Cost model decision regions]{Decision regions of Equation~\ref{eq:cost} for the pandas path
    ($\rho = 1$, $R = R' = 1$\,GFLOP/s) at three nodes on a 500\,MB/s link.
    The step at 512\,MB is a separate rule, not Equation~\ref{eq:cost}: work
    below 8 flop per byte stays local up to 512\,MB, so grouped aggregation and
    joins are distributed only above it.}
    \label{fig:cost}
\end{figure}

The consequence for dense matrix multiplication is that it stays local on a CPU
cluster: local BLAS at about 200\,GFLOP/s outruns what a CPU cluster's transport
can feed it at practical sizes. Measured against a four-core peer, local won by
25 to 100 times across four shapes, and the model chose local every time. A
product ships when it does not fit locally, when a GPU peer makes the model
favour it, or when forced. Singular value decomposition, two orders of magnitude
slower locally, does ship.

% ------------------------------------------------------------
\section{Fault Tolerance}
\label{sec:ft}

The invariant is that using the cluster is never worse than not using it. Five
independent layers enforce it for parallel work; a training ring that loses a
rank continues on the survivors (Section~\ref{sec:ddp}).

\begin{enumerate}[leftmargin=*, label=\arabic*.]
  \item \textbf{The race.} Chunks are dealt alternately to the local side and the
  cluster (Figure~\ref{fig:race}), and results fill a slot table first-wins. Idle local cores re-run any
  chunk the cluster has not yet returned, and the call returns only once every
  item is covered, so a dead cluster degrades to a plain local pool.
  \item \textbf{Per-primitive fallback.} Every intercepted call is wrapped, and
  any exception falls back to the original implementation; the one exception is
  a matrix product that fits neither locally nor on the cluster, which raises.
  \item \textbf{Peer fallback.} A failed peer is skipped for the rest of that
  chunk and the next candidate tried; if all fail, the work runs locally. A slow
  peer is raced by a copy of its part.
  \item \textbf{Worker respawn.} A dead warm worker is restarted, falling back to
  a one-shot process if that fails.
  \item \textbf{Rescheduling.} A failed job is placed again with back-off, on an
  untried node first and at most three times (Section~\ref{sec:lease}).
\end{enumerate}

\begin{figure}[!htbp]
    \centering
    \includegraphics[width=150mm]{\figpath{fig_4.12_race_timeline.png}}
    \caption[Racing local and remote work]{Racing local and remote work. The local side is always able to
    finish alone, which is what makes the guarantee unconditional.}
    \label{fig:race}
\end{figure}

Memory is bounded at the same time: a parallel-map request larger than a node
can grant is forwarded to a peer that can, one exceeding 40\% of that memory is
split into sequential parts, and a request body is capped at the smaller of
2\,GiB and half the daemon's memory budget.

% ------------------------------------------------------------
\section{Resource Estimation and Run Reporting}
\label{sec:observe}

\textbf{Resource estimation.} Tiers are consulted in order, the first that
yields a value winning: an explicit \texttt{-{}-mem}; 1.2 times the largest peak
memory of the last five runs of the same script; or per-library baselines plus
100\,MB for the interpreter. The default when nothing is known is 512\,MB. The
historical tier closes a feedback loop: the worker samples the job's process tree
while it runs and records the peak, which sizes the next run's admission request.

\textbf{Commands.} The CLI has thirteen commands, among them setup, invite,
connect, run, status, attach, doctor and uninstall. Configuration is a single file
holding the secret, the introducers and the node's role, and the ports are one
block: 4242 for the daemon API and discovery, 4243 for direct links and 4244 for
the introducer. \texttt{pipedpeer doctor} checks the daemon, the secret, each
introducer and each peer (TCP and UDP separately), the firewall, user namespaces,
\texttt{crun}, Nix, the GPU driver, disk and clock, each with its fix. A plain run
behaves like Python: it prints the script's own output, passes the script its
arguments and returns its exit code, adding one line on where the work ran when
any left the machine.

\textbf{The run record.} Every run keeps one record of timed steps ---
serialising, sending, queueing, computing and receiving each chunk on each
machine, building and sending the environment, each synchronisation round, and
the script itself --- written by the shim, the daemon and the client, and shown
by \texttt{-v}, the job history, the web page opened by \texttt{pipedpeer ui},
and as CSV or JSON. \texttt{pipedpeer events} prints in plain language what this
machine's daemon has done, and a dry run reports, before anything runs, the
memory estimate, unresolved imports, the parallel calls that can spread, and
where the job would go.
